% p2-10-work-packages

\section{P2 §10. Work Packages, Falsifiability, and Reviewer Risk}

10. Work Packages, Falsifiability, and Reviewer Risk

      Figure (fig-p2-ch10-work-packages). Work-packages triptych --- WP1..WP5 grid / Falsifiability
                                    arrows / Reviewer-risk shield.

  The programme is organized into five work packages with explicit success and failure
  criteria. The complete formal success/failure table is given in Section 4.4; this section
  provides operational detail and confronts the objections a referee would most likely
  raise.

  WP1: Search for $\varphi$-Eigenvalue Ratios at CFT Fixed Points
  Task. Enumerate known unitary 2D CFTs (Virasoro minimal models, N=2 models, WZW
  models at low level) and check whether any have two primary operator dimensions
  Delta1, Delta2 with Delta2/Delta1 = $\varphi$n for small integer n.

  Method. Direct computation from the Kac table and known character formulas. This is a
  finite computation for each minimal model. The c < 1 Virasoro minimal models have
  (p,p') pairs with p, p' $\leq$ 30 for the first dozen models; the dimension ratios are rational,
  and checking Deltar',s'/Deltar,s $\in$ Q(5) is a decidable algebraic question.

  Success criterion. At least one unitary, modular-invariant CFT partition function
  contains a $\varphi$-eigenvalue pair (see Section 4.4 table).

  Failure criterion. Exhaustive check over all c < 1 minimal models plus the first five
  WZW levels at low rank finds no such pair.

  Epistemic status if successful. Establishes existence of $\varphi$-towers in consistent CFTs;
  does not establish relevance to Standard Model or to Paper 1 data. Upgrade from L5 to
  L3.

  Epistemic status if failed. Strongly suggests $\varphi$-towers are not generic within the known
  landscape of simple unitary CFTs, weakening the theoretical motivation for the
  mechanism programme. Programme pivots to WP3 or terminates.

  Timeline estimate. 1--3 months for a computational algebraist or conformal field theorist
  with access to standard software (Mathematica, Kac-table enumeration code).

  WP2: One-Loop Toda Mass Spectrum Under $\varphi$-DSI
  Task. For the H3 affine Toda field theory of Fring and Korff (2005), compute quantum
  (one-loop) corrections to the mass spectrum and determine whether the classical
  $\varphi$-ratio survives renormalization. Note that the Fring-Korff paper already determines
  one-loop corrections to mass renormalization (arXiv:hep-th/0506226); this work
  package is to extract the explicit ratio from those results and determine its deviation
  from $\varphi$.

  Method. Standard one-loop correction using the known coupling structure, following
  the procedure of Fring and Korff.

  Success/failure criterion. See Section 4.4 table, row 3. The key number is |delta m2/m2
  - delta m1/m1|, i.e., the difference in relative mass corrections; if this is parametrically
  small in the coupling beta, the classical ratio is perturbatively stable.

  Timeline estimate. 1--2 months for a specialist in Toda field theory or integrable models.

  WP3: DSI Ratio Derivation from First Principles

  Task. Identify whether any consistent UV completion (string-theory compactification,
  holographic dual, or lattice gauge theory) constrains an RG flow to discrete dilations
  with ratio $\varphi$ rather than a generic irrational.

  Method. Literature survey of holographic RG flows with discrete structure; candidate
  Calabi-Yau compactifications with H4-type symmetry in their Kähler moduli space;
  icosahedral quasicrystal effective field theory as a condensed matter analogue. A key
  question is whether any orbifold compactification with icosahedral holonomy gives a
  DSI ratio constrained by the H3 dilation spectrum.

  Success/failure criterion. See Section 4.4 table, row 4.

  Timeline estimate. 6--18 months; this is the most open-ended work package and
  constitutes a genuine research problem.

  WP4: Precision Test of Log-Periodic Running
  Task. Given the log-periodic correction template of equation (21), determine the
  sensitivity required to detect DSI amplitudes Ai of order 10-3 in running coupling
  measurements, and compute the current upper bound from LEP/LHC data.

  Method. Propagate LEP/LHC coupling measurement uncertainties through a chi2 test
  for log-periodic deviations with period fixed at ln$\varphi$ approx 0.4812.

  Current bound estimate. From the precision of alphas(mZ) measurements (approx 0.3\%
  at LEP) and of electroweak coupling unification checks, the absence of log-periodic
  oscillations with period ln$\varphi$ over the range mu $\in$ [mZ, 1 TeV] constrains Ai $\leq$sssim few
  $\times$ 10-3. The computation must determine whether this bound is already strong enough
  to exclude the physically motivated range or whether FCC-ee precision
  (deltaalphas/alphas $\sim$ 0.03\%) is needed.

  Success/failure criterion. See Section 4.4 table, row 6.

  Timeline estimate. 2--4 months given access to the LEP combination data and standard
  coupling-running codes.

  WP5: Robustness Audit of Paper 1 Compressibility
  Task. Systematically test whether the symbolic compressibility result of Paper 1
  survives: (a) changes in numerical tolerance epsilon; (b) alternative encoding bases; (c)
  sampling of different coupling combinations including from the flos34 catalogue; (d)
  comparison against degree-2 algebraic constants of similar magnitude (2, 3, 6).

  Success/failure criterion. See Section 4.4 table, row 7.

  Note on flos34. The 42-parameter catalogue from the Flos Aureus monograph is used
  here as candidate input, not as independent confirmation. Any entry from flos34 used in
  WP5 must be treated as an unlabeled data point subject to the same statistical tests as
  any other coupling combination. To avoid circular risk, the enumeration of flos34
  entries should be conducted by a party who has not seen the Paper 1 compressibility
  scores.

  Timeline estimate. 2--3 months for the full audit; 1 month for the baseline
  tolerance-variation test.

  10.1 Reviewer Risk Assessment and Responses
  The following table anticipates the objections most likely to be raised by a specialist
  referee and provides the paper's response.

\subsection{References}

- Coldea, R. et al. "Quantum Criticality in an Ising Chain: Experimental Evidence for Emergent E8 Symmetry." \textit{Science} 327 (2010) 177. DOI \href{https://doi.org/10.1126/science.1180085}{10.1126/science.1180085}.
- Wu, J. et al. "E8 Symmetry in CoN$b_{2}$$O_{6}$ at Higher Resolution." \textit{Phys. Rev. B} 108, L060403 (2023). DOI \href{https://doi.org/10.1103/PhysRevB.108.L060403}{10.1103/PhysRevB.108.L060403}.
- Fring, A.; Korff, C. "Affine Toda field theories related to Coxeter groups of non-crystallographic type." \textit{Nucl. Phys. B} 729 (2005) 361. \href{https://arxiv.org/abs/hep-th/0506226}{arXiv:hep-th/0506226}. DOI \href{https://doi.org/10.1016/j.nuclphysb.2005.08.044}{10.1016/j.nuclphysb.2005.08.044}.
- Pellis, S. "The Fine-Structure Constant and the Golden Ratio." \href{https://vixra.org/pdf/2110.0117v4.pdf}{viXra 2110.0117 v4}.
